\section{Useful Properties related to SL(2, R) covariance}
\label{sec:sl2r_covariance}

\subsection{Modular forms}
For a function $F \to |c\tau+d|^k F$, we can construct a covariant derivative
\begin{equation}
\hat{\partial}_m F \equiv \left[\partial_m - \frac{k}{2}\partial_m\Phi \right] F,
\label{U(1)-covariant derivative}
\end{equation}
which can equivalently be written as
\begin{equation}
\hat{\partial}_m F = e^{k\phi/2}\,\partial_m \left(e^{-k\phi/2} F\right),\quad
\phi\equiv \Phi - \Phi_0,
\label{SL2R_completion_by_scaling}
\end{equation}
for some constant $\Phi_0$ that can be understood as the constant part of $\Phi$. If the symmetry group were $SL(2, \mathbb{R})$, then one may set $\Phi_0 = 0$ via an appropriate $SL(2, \mathbb{R})$ transformation. However, this is not the case for $SL(2, \mathbb{Z})$, so we will keep $\Phi_0$ explicit and conventionally define $g_s \equiv e^{\Phi_0}$.

\subsection{IIB fields under \texorpdfstring{$SL(2, \mathbb{R})$}{SL(2, R)}}
For an $SL(2, \mathbb{R})$ transformation
\begin{equation}
\begin{pmatrix}
a & b\\
c & d
\end{pmatrix}
\in SL(2, \mathbb{R}),
\end{equation}
the axio-dilaton, defined as
\begin{equation}
\tau \equiv C + ie^{-\Phi},
\end{equation}
transforms as
\begin{equation}
\tau \to \frac{a\tau + b}{c\tau + d}.
\end{equation}
The dilaton and axion individually transform as
\begin{equation}
e^{-\Phi} \to \frac{e^{-\Phi}}{|c\tau + d|^2},\qquad
C \to \frac{ac|\tau|^2 + (ad+bc)C +bd}{|c\tau+ d|^2}.
\end{equation}
The modulud $|\tau|^2$ transforms as
\begin{equation}
|\tau|^2 \to \frac{|a\tau + b|^2}{|c\tau + d|^2} = \frac{a^2|\tau|^2 + 2abC + b^2}{|c\tau + d|^2}.
\end{equation}
The string frame metric $G_{mn}$ transforms as
\begin{equation}
G_{mn}\to |c\tau + d|\, G_{mn},\qquad
\sqrt{-G} \to |c\tau + d|^5 \sqrt{-G}.
\label{eq:metric-sl2z}
\end{equation}

\subsection{\texorpdfstring{$SL(2, \mathbb{R})$}{SL(2, R)} completion}
From \eqref{eq:metric-sl2z}, the string frame metric carries $U(1)$ charge $k=1$ under $SL(2, \mathbb{R})$.
Since the Ricci scalar $R$ contains derivatives acting on $G_{mn}$, it will not transform covariantly under $SL(2, \mathbb{R})$.
However, we may still consider its $SL(2, \mathbb{R})$ completion $\hat{R}$ obtained by promoting every ordinary derivative inside $R$ to the $SL(2, \mathbb{R})$-covariant derivative defined in \eqref{U(1)-covariant derivative}.
One finds
\begin{equation}
\hat{R}(G_{mn},\Phi) = e^{-\phi/2}R(e^{-\phi/2}G_{mn}),\quad
\hat{R} \to |c\tau + d|^{-1}\hat{R},
\end{equation}
and
\begin{equation}
\sqrt{-G}e^{-2\Phi}\hat{R}(G_{mn},\Phi)
=\sqrt{-G}e^{-2\Phi}\left[R(G_{mn})+\frac{9}{2}(\partial\Phi)^2\right]+\partial(...).
\end{equation}

% \section{Reduction on torus}
% The compactification of a $D$-dimensional metric on a torus $T^2$ is given by
% \begin{equation}
% ds^2_D = e^{2\alpha\varphi}g_{mn}dx^m dx^n + e^{2\beta\varphi}G_{ij}(dy^i + A^i_m dx^m)(dy^j + A^j_n dx^n)
% \end{equation}
% The Ricci tensor is given by
